% Generated draft: requires mathematical and source review.
\documentclass[11pt]{article}
\usepackage[margin=25mm]{geometry}
\usepackage{fontspec}
\setmainfont{Latin Modern Roman}
\usepackage{amsmath,amssymb,mathtools,unicode-math}
\setmathfont{Latin Modern Math}
\usepackage{enumitem,needspace,xurl,hyperref}
\setlength{\parindent}{0pt}
\setlength{\parskip}{5pt}
\setlength{\emergencystretch}{3em}
\newcommand{\N}{\mathbb{N}}
\newcommand{\Z}{\mathbb{Z}}
\newcommand{\Q}{\mathbb{Q}}
\newcommand{\R}{\mathbb{R}}
\newcommand{\C}{\mathbb{C}}
\newcommand{\F}{\mathbb{F}}
\newcommand{\A}{\mathbb{A}}
\newcommand{\PP}{\mathbb{P}}
\newcommand{\E}{\mathbb{E}}
\newcommand{\eps}{\varepsilon}
\newcommand{\dd}{\,\mathrm{d}}
\begin{document}
\section*{UnsolvedMath ID: 30002539}
\textbf{Source problem code:} catalogue entry
\subsection{Definitions and mathematical statement}
The supplied record does not define the expressions ``the Hopf case'' or ``the graphical calculus for operations in Deligne's conjecture.'' In particular, it specifies neither the Hopf objects under consideration nor the operations, graphs, relations, or interpretation maps that an analogue should describe. These choices remain unresolved; no particular algebraic setting is imposed here.

The source poses the following construction question:
\begin{quote}
What is the analogue, for the Hopf case, of the graphical calculus for operations in Deligne's conjecture?
\end{quote}
A precise construction problem requires identification of the intended Hopf setting, the particular Deligne-type graphical calculus to be generalized, and the properties that the proposed analogue must satisfy. The supplied record does not provide these specifications, so it does not determine a uniquely formulated mathematical assertion.
\subsection{Short English statement}
The problem asks for a Hopf analogue of a graphical calculus for operations associated with Deligne's conjecture. The intended Hopf setting and the reference calculus are unspecified, so the target construction requires clarification from the source context.
\subsection{Sources}
\begin{itemize}
\item Supplied local UnsolvedMath registry record, permanent ID \(30002539\), source problem code \texttt{catalogue entry}, titled ``Hopf Graphical Calculus for Deligne-Type Operations,'' in the Oberwolfach Reports Open Problems collection. The record lists proposed year 2014 and identifies no individual proposer.
\item \emph{Mini-Workshop: Batalin-Vilkovisky Algebras, Operads, and Hopf Algebroids} (2014), Oberwolfach Reports, pp.~595--634, DOI 10.4171/owr/2014/11. Supplied source report and literature reference labeled ``Oberwolfach Report 11, 2014'': \url{https://doi.org/10.4171/owr/2014/11}.
\item Supplied literature reference labeled ``Kaufmann and Ward, 2014'': \url{https://arxiv.org/abs/1312.1269}.
\item Literature assessment dated exactly \texttt{2026-08-21}: the supplied curation labels the problem open and states that Feynman categories provide Hopf-algebra constructions, but that no general analogue of the cited Deligne-type graphical calculus for the Hopf case has been established. This is supplied dated curation derived from the Oberwolfach Reports material, not an independently verified assessment; its status and cited claims require independent verification.
\end{itemize}
\end{document}
